% !Mode:: "TeX:UTF-8"

\section{项目设计目标}

随着物联网技术和智能家居概念的普及，传统家电的智能化改造已成为当前技术发展的重要方向。电风扇作为家庭常见电器，其控制方式仍以机械开关和物理调速为主，存在操控不便、功能单一等问题。基于嵌入式微控制器的智能风扇系统，能够通过先进的电子控制技术实现更人性化、智能化的操作体验。

STM32系列微控制器凭借其强大的处理能力、丰富的外设接口（如定时器、PWM、ADC、多种通信接口）以及低功耗特性，为家电智能化提供了理想的解决方案。本项目结合嵌入式系统课程所学知识，选择STM32作为核心控制器，设计并实现一款可通过红外遥控操作的智能电风扇系统，旨在将理论知识与实际应用相结合，探索嵌入式系统在家电控制领域的应用潜力。

本项目旨在设计并实现一个功能完备、操作便捷的遥控电风扇系统，主要目标包括：
\begin{itemize}
    \item \textbf{实现电机精准控制：} 利用STM32的定时器模块产生可调PWM信号，驱动直流电机实现多档位或无级调速功能。
    \item \textbf{开发红外遥控交互：} 设计红外接收电路，实现遥控器信号的解码与响应，支持开关机、风速调节等基本操作。
    \item \textbf{增加定时功能：} 实现用户可设定的定时关闭功能（如1小时、2小时、4小时等选项），提升产品实用性。
    \item \textbf{完成系统集成：} 设计完整的硬件电路（含电机驱动模块、红外接收模块、电源模块等），并开发相应的嵌入式软件系统。
    \item \textbf{实现原型验证：} 在STM32开发板（如STM32F103C8T6最小系统板）及外围电路搭建的硬件平台上进行功能测试与性能评估。
\end{itemize}
